% =====================================================================
% PLANTILLA: Guía Didáctica de Aprendizaje - Colegio San Alberto Magno
% =====================================================================
% CÓMO USAR ESTA PLANTILLA
% 1. Copia este archivo con un nombre nuevo, p. ej. guia-07-fracciones.tex
% 2. Deja guia-didactica.sty y logo-colegio.png en la misma carpeta
%    (o en una carpeta compartida referenciada con \graphicspath).
% 3. Completa los datos de la sección "DATOS DE LA GUÍA" más abajo.
% 4. Compila con:  pdflatex guia-07-fracciones.tex
% =====================================================================
\documentclass[11pt]{article}
\usepackage{guia-didactica}

% =====================================================================
% DATOS DE LA GUÍA  (edita solo esta sección para una guía nueva)
% =====================================================================
\newcommand{\NumeroGuia}{1}
\newcommand{\Area}{Matemáticas}
\newcommand{\Asignatura}{Matemáticas}
\newcommand{\Grado}{7°}
\newcommand{\Periodo}{III}
\newcommand{\Docente}{Nombre del Docente}
\newcommand{\Estudiante}{\hrulefill}
\newcommand{\FechaInicio}{}
\newcommand{\FechaEntrega}{}

% Proyectos transversales articulados: cambia 'x' por vacío para des-marcar
\newcommand{\MarcaEDPC}{}
\newcommand{\MarcaPEEF}{}
\newcommand{\MarcaEVS}{}
\newcommand{\MarcaPRAE}{}
\newcommand{\MarcaCOLOM}{}
\newcommand{\MarcaNOCEST}{}
\newcommand{\MarcaOTRO}{}

\begin{document}

\GuiaTitulo{\NumeroGuia}

% Renglón de proyectos transversales del título (como en el original)
\begin{center}
\small
Educación para la Democracia, Paz y Convivencia (EDPC) \hspace{1cm}
P. Ambiental Escolar (PRAE)\\
Educación Económica y Financiera (PEEF) \hspace{1cm}
Estilos de Vida Saludable (EVS)
\end{center}

% ---------------------------------------------------------------------
% Datos básicos
% ---------------------------------------------------------------------
\begin{gdDatos}
\textbf{Área:} & \Area & \textbf{Grado:} & \Grado \\
\textbf{Asignatura:} & \Asignatura & \textbf{Período:} & \Periodo \\
\textbf{Docente:} & \Docente & \textbf{Fecha de inicio:} & \FechaInicio \\
\textbf{Estudiante:} & \Estudiante & \textbf{Fecha de entrega:} & \FechaEntrega \\
\end{gdDatos}

% ---------------------------------------------------------------------
% 1. Identificación y Propósito de Aprendizaje
% ---------------------------------------------------------------------
\gdSeccion{Identificación y Propósito de Aprendizaje}

\gdCampo{Eje Temático:}{}
\gdRenglones{1}

\gdCampo{DBA:}{}
\gdRenglones{2}

\gdSubseccion{Proyecto Transversal Articulado}
\noindent
\gdCasilla{EDPC}{\MarcaEDPC}\hfill
\gdCasilla{PEEF}{\MarcaPEEF}\hfill
\gdCasilla{EVS}{\MarcaEVS}\hfill
\gdCasilla{PRAE}{\MarcaPRAE}\hfill
\gdCasilla{COLOM}{\MarcaCOLOM}\hfill
\gdCasilla{NOCEST}{\MarcaNOCEST}\hfill
\gdCasilla{OTRO}{\MarcaOTRO}

\vspace{2mm}
\gdCampo{Pregunta Problematizadora del contexto:}{}
\gdRenglones{2}

\gdCampo{Propósito Holístico de la Guía:}{}
\gdRenglones{2}

% ---------------------------------------------------------------------
% 2. Ruta de Aprendizaje y Secuencia Didáctica
% ---------------------------------------------------------------------
\gdSeccion{Ruta de Aprendizaje y Secuencia Didáctica}

\gdSubseccion{Momento No. 1 -- Exploración}
\gdCampo{Activación de saberes y conexión emocional:}{}
\gdRenglones{1}
\gdCampo{¿Qué sabemos del tema?:}{}
\gdRenglones{1}
\gdCampo{Reto de apertura:}{}
\gdRenglones{1}
\gdCampo{Diálogo de valores:}{}
\gdRenglones{1}
\gdCampo{Adaptación DUA / PIAR para la Exploración:}{}
\gdRenglones{1}

\gdSubseccion{Momento No. 2 -- Estructuración (Construyo mi conocimiento)}
\gdCampo{Etapa Concreta (Saber con las manos):}{}
\gdRenglones{1}
\gdCampo{Etapa Pictórica (Saber con imágenes):}{}
\gdRenglones{1}
\gdCampo{Etapa Abstracta (Saber con símbolos):}{}
\gdRenglones{1}
\gdCampo{Ajuste Razonable (PIAR):}{}
\gdRenglones{1}

\gdSubseccion{Momento No. 3 -- Práctica y Transferencia (Aplico y transformo mi entorno)}
\gdCampo{Trabajo Colaborativo / Tutoría de Pares:}{}
\gdRenglones{1}
\gdCampo{Misión Transversal (PRAE -- ABP):}{}
\gdRenglones{1}
\gdCampo{Demostración de lo aprendido:}{}
\gdRenglones{1}

\gdSubseccion{Momento No. 4 -- Valoración y Cierre (Reflexiono sobre mi proceso)}
\gdCampo{El error como aprendizaje:}{}
\gdRenglones{1}
\gdCampo{Aportes a mi comunidad:}{}
\gdRenglones{1}

% ---------------------------------------------------------------------
% 3. Matriz de Evaluación Formativa y Cualitativa
% ---------------------------------------------------------------------
\gdSeccion{Matriz de Evaluación Formativa y Cualitativa}

\begin{gdMatrizEvaluacion}
Cognitiva (Saber) & & $\square$ & $\square$ & $\square$ \\
Procedimental (Hacer) & & $\square$ & $\square$ & $\square$ \\
Socioafectiva (Ser) & & $\square$ & $\square$ & $\square$ \\
Trascendente (Convivir) & & $\square$ & $\square$ & $\square$ \\
\end{gdMatrizEvaluacion}

% ---------------------------------------------------------------------
% 4. Autoevaluación y Reflexión Metacognitiva
% ---------------------------------------------------------------------
\gdSeccion{Autoevaluación y Reflexión Metacognitiva}

\gdCampo{Lo que más me gustó de esta guía fue:}{}
\gdRenglones{1}
\gdCampo{La mayor dificultad que tuve y cómo la superé:}{}
\gdRenglones{1}
\gdCampo{Mi compromiso para la próxima guía es:}{}
\gdRenglones{1}

\vspace{4mm}
\noindent\textbf{Firma del Estudiante:} \hrulefill \hspace{1cm}
\textbf{Firma Docente:} \hrulefill

% ---------------------------------------------------------------------
% 5. Seguimiento y Apoyos (Diligenciado por el docente)
% ---------------------------------------------------------------------
\gdSeccion{Seguimiento y Apoyos (Diligenciado por el docente)}

\noindent\textbf{Nivel de desempeño alcanzado:}\quad
$\square$ Inicial \quad $\square$ En Proceso \quad $\square$ Consolidado

\vspace{2mm}
\gdCampo{Ajustes DUA/PIAR aplicados en la guía:}{}
\gdRenglones{2}

\gdCampo{Observaciones y recomendaciones pedagógicas:}{}
\gdRenglones{3}

% ---------------------------------------------------------------------
% 6. Seguimiento Académico (Diario de Campo)
% ---------------------------------------------------------------------
\gdSeccion{Seguimiento Académico (Diario de Campo)}

\noindent\begin{tabularx}{\linewidth}{|X|X|l|X|}
\hline
\textbf{Semana No.} & & \textbf{Fecha:} & \\
\hline
\textbf{No. de estudiantes:} & & \textbf{No. de estudiantes ausentes:} & \\
\hline
\multicolumn{4}{|l|}{\textbf{Faltan:}\vphantom{Ay}} \\
\hline
\multicolumn{4}{|l|}{\textbf{Tema:}\vphantom{Ay}} \\
\hline
\multicolumn{4}{|l|}{\textbf{Observación:}\vphantom{Ay}} \\
\hline
\end{tabularx}

\end{document}
